% Figure for Classical Mechanics §A3.10 (Theorem: falling from rest with linear drag; v(t) approaches terminal speed; compared with free fall).
\begin{tikzpicture}
  \begin{axis}[nfig, width=9cm, height=5.6cm, xmin=0, xmax=5.4, ymin=0, ymax=1.35,
      xlabel={$t/(m/b)$}, ylabel={$v/v_T$},
      xtick={0,1,2,3,4,5}, ytick={0,0.632,1}, yticklabels={0,0.63,1},
      clip=false]
    \addplot[black!45, thin, dashed, domain=0:5.4] {1};
    \addplot[black!35, thin, densely dotted] coordinates {(1,0) (1,0.632) (0,0.632)};
    \addplot[figB, thick, dashed, domain=0:1.3] {x};
    \addplot[figR, very thick, domain=0:5.4, samples=80] {1-exp(-x)};
    \node[font=\scriptsize, figB, anchor=west] at (axis cs:1.32,1.28) {free fall, $v = gt$};
    \node[font=\scriptsize, figR, anchor=north west] at (axis cs:2.6,0.9) {$v = v_T\,(1 - e^{-bt/m})$};
    \node[font=\scriptsize, black!60, anchor=south east] at (axis cs:5.4,1.0) {terminal speed $v_T = mg/b$};
  \end{axis}
\end{tikzpicture}
